\fancyhead[L]{\small Contributions to the Theory}
\SourceLeaf{080}
\begin{center}
{\large\bfseries CONTRIBUTIONS TO THE THEORY\par}
\medskip
{\Large\bfseries OF ALGEBRAIC EQUATIONS.\par}
\medskip
\rule{0.18\linewidth}{0.4pt}
\end{center}

This memoir treats two different subjects concerning algebraic equations. First, I present the proof of the fundamental theorem of the theory of algebraic equations that I gave fifty years ago, in a modified form and with substantial additions. The second part is devoted to a special treatment of algebraic equations with three terms and contains methods for easily determining not only their real but also their imaginary roots.

\begin{center}
\rule{0.18\linewidth}{0.4pt}\par
\medskip
\textbf{FIRST DIVISION.}
\end{center}

The memoir published in 1799, \emph{A New Proof of the Theorem That Every Polynomial in One Variable Can Be Resolved into Real Factors of the First or Second Degree}, had two aims: first, to show that all the proofs previously attempted of this most important theorem of the theory of algebraic equations were insufficient and illusory, and second, to give a new, fully rigorous proof. There is no need to return to the first subject. I myself later followed the new proof given there with two others, and a fourth was first given by \textsc{Cauchy}. These four proofs \SourceLeaf{081}all rest on as many different foundations. They agree, however, in that each initially proves only the existence of \emph{one} factor of the function in question. This certainly does not detract from the rigour of the proofs. For it is clear that removing this one factor from the given function leaves a similar function of lower degree, to which the theorem can be applied again, and that repeating the procedure must ultimately give a complete factorization of the original function into factors of the specified kind. Nevertheless, every proof undoubtedly attains a higher degree of completeness if it can be shown to make the existence of all the factors directly evident. I already indicated in that memoir, Article 23, that the first proof does this, without developing the point further there. This shall now be supplied, and I take the opportunity to repeat the main points of the entire proof in a modified form that, I believe, offers greater clarity. As for the outward formulation of the theorem itself, the form used in 1799, that the function $x^n+Ax^{n-1}+Bx^{n-2}+\mathrm{etc.}$ can be resolved into real factors of the first or second degree, was chosen then because all use of imaginary quantities was to be avoided. Now that the concept of complex quantities is familiar to everyone, it seems more appropriate to abandon that form and state the theorem by saying that the function can be resolved into $n$ \emph{linear} factors. The constant parts of these factors need not be real quantities; any complex values must also be allowed for them. With this formulation the theorem itself gains in generality, since the restriction to real values need not be imposed on the coefficients $A,B$ etc. either; any values remain admissible for them.

\FifthArticle{1}
We therefore consider the function of the indeterminate $x$
\[
x^n+Ax^{n-1}+Bx^{n-2}+\mathrm{etc.}+Mx+N=X
\]
where $A,B\,.\,.\,.\,.\,M,N$ denote fixed real or imaginary coefficients.
\par
\SourceLeaf{082}
The relation between the roots of the equation $X=0$ and the linear factors of $X$ is known from elementary algebra. Namely, if the substitution $x=p$ satisfies that equation, then $x-p$ is a factor of $X$. If there are $n$ different ways of satisfying it, namely $x=p$, $x=p'$, $x=p''$ etc., then the product $(x-p)$ $(x-p')(x-p'')\, .\, .\, .\, .$ is identical to $X$. Under special circumstances, however, a solution such as $x=p$ may also give the factor $(x-p)^2$, or $(x-p)^3$, or some higher power in $X$. In these cases the root $p$ is counted as occurring twice, three times, and so on.

Thus, if only a proof that the function $X$ certainly admits \textit{one} linear factor is required, it suffices to establish the existence of any root of the equation $X=0$. But if the complete factorization of the function into linear factors is to be proved all at once, it must be shown that the equation $X=0$ can be satisfied either by $n$ unequal values of $x$, or by a smaller number of unequal solutions, some of which have the characteristics of equal roots counted more than once, in such a way that counting all the unequal and equal roots gives the total $=n$.

